\documentclass[12pt]{article}
\usepackage{amsmath}
\usepackage{amssymb}
\usepackage{physics}
\usepackage{graphicx}
\usepackage{hyperref}

\newcommand{\ds}{\displaystyle}

\title{The Foucault Pendulum}
\author{}
\date{}

\begin{document}
\maketitle

\section{Introduction}

The equations needed below are obtained in the preceding essay \href{/physics/classical-mechanics/rigid-body-motion/motion-on-the-rotating-earth}{Motion on the Rotating Earth} by linearizing
the equations of motion near a point on the Earth's surface. Defining dimensional
variables of length as  

\begin{align}
  x &= \rho_0\,\phi & y &= \rho_0\sin(\phi_0)\,\theta & z &= \rho
\end{align}

and defining a local coordinate system as

\begin{align}
  \vb{e}_z &= \vb{e}_{\rho_0} &
  \vb{e}_y &= \vb{e}_{\theta_0} &
  \vb{e}_x &= \vb{e}_{\phi_0}
\end{align}

the linearized equations of motion are

\begin{equation}
  \begin{aligned}
    \ddot{z}
    &=
    -g
    +\omega^2\rho_0\sin^2\phi_0
    +\omega^2\sin^2\phi_0\,z
    +\omega^2\sin(2\phi_0)\,x
    +2\omega\sin\phi_0\,\dot{y},
    \\
    \ddot{x}
    &=
    \frac{1}{2}\omega^2\rho_0\sin(2\phi_0)
    +\frac{1}{2}\omega^2\sin(2\phi_0)\,z
    +\omega^2\cos(2\phi_0)\,x
    +2\omega\cos\phi_0\,\dot{y},
    \\
    \ddot{y}
    &=
    -2\omega\sin\phi_0\,\dot{z}
    -2\omega\cos\phi_0\,\dot{x}.
  \end{aligned}
  \label{earth:local-cartesian-equations}
\end{equation}

\section{The Foucault Pendulum}

In this section, we derive the equations of motion of the Foucault pendulum by
applying several simplifying assumptions. The Foucault pendulum is one of the
most elegant experiments demonstrating the rotation of the Earth. It consists of
a heavy mass $m$ suspended by a long wire or rigid rod of length $l$, free to
oscillate in any vertical plane. As the pendulum oscillates, the plane of motion
slowly rotates with respect to the Earth due to the Coriolis force.

Before proceeding, it is useful to estimate the relative magnitude of the terms
appearing in the equations of motion. The angular velocity of the Earth is

\[
  \omega
  =
  \frac{2\pi}{T_{\mathrm{sid}}}
  =
  \frac{2\pi}{86\,164\ \mathrm{s}}
  \approx
  7.2921\times10^{-5}\ \mathrm{s}^{-1},
\]

where $T_{\mathrm{sid}}$ denotes one sidereal day.

For a typical Foucault pendulum with length between $15\,\mathrm{m}$ and
$30\,\mathrm{m}$, the natural angular frequency is

\[
  \omega_0
  =
  \sqrt{\frac{g}{l}}
  \approx
  0.57\text{--}0.81\ \mathrm{s}^{-1},
\]

and therefore

\[
  \frac{\omega}{\omega_0}
  =
  \omega\sqrt{\frac{l}{g}}
  \sim
  10^{-4}.
\]

The Coriolis terms in the equations of motion are proportional to $\omega$,
whereas the centrifugal corrections are proportional to $\omega^2$. We shall
therefore retain only terms up to first order in the Earth's angular velocity,
neglecting all terms of order $\omega^2$. Equation
\eqref{earth:local-cartesian-equations} consequently reduces to

\begin{equation}
  \begin{aligned}
    \ddot{x}
      &=
      2\omega\cos(\phi_0)\,\dot{y},
      \\
    \ddot{y}
      &=
      -2\omega\sin(\phi_0)\,\dot{z}
      -2\omega\cos(\phi_0)\,\dot{x},
      \\
    \ddot{z}
      &=
      -g
      +2\omega\sin(\phi_0)\,\dot{y}.
  \end{aligned}
  \label{earth:first-order-local-equations}
\end{equation}

These equations describe the motion of a particle subjected only to gravity and
the Coriolis force. The constraint force exerted by the pendulum wire will be
introduced in the next section.

Introducing the position vector

\[
  \vb{p}
  =
  x\,\vb{e}_x
  +
  y\,\vb{e}_y
  +
  z\,\vb{e}_z,
\]

the Coriolis force can be written in vector form as

\[
  \vb{F}_{\mathrm{cor}}
  =
  -2m\,
  \boldsymbol{\omega}
  \times
  \dot{\vb{p}},
\]

where the Earth's angular velocity expressed in the local Cartesian basis
introduced in the previous section is

\[
  \boldsymbol{\omega}
  =
  \omega
  \left(
      -\sin(\phi_0)\,\vb{e}_x
      +
      \cos(\phi_0)\,\vb{e}_z
  \right).
\]

Explicitly,

\[
  \vb{F}_{\mathrm{cor}}
  =
  m
  \begin{pmatrix}
      2\omega\cos(\phi_0)\,\dot{y}
      \\
      -2\omega\sin(\phi_0)\,\dot{z}
      -
      2\omega\cos(\phi_0)\,\dot{x}
      \\
      2\omega\sin(\phi_0)\,\dot{y}
  \end{pmatrix}.
\]

The Coriolis force admits the generalized potential

\[
  V_{\mathrm{cor}}
  =
  -m
  \,
  \ip{\boldsymbol{\omega}\times\vb{p}}
     {\dot{\vb{p}}},
\]

which, in local Cartesian coordinates, becomes

\[
  V_{\mathrm{cor}}
  =
  m\omega
  \left(
      y\cos(\phi_0)\dot{x}
      -
      \bigl(
          \sin(\phi_0)z
          +
          \cos(\phi_0)x
      \bigr)
      \dot{y}
      +
      \sin(\phi_0)y\dot{z}
  \right).
\]

Including gravity, the corresponding Lagrangian is therefore

\[
  \begin{aligned}
    L
    &=
    \frac12m
    \left(
        \dot{x}^{\,2}
        +
        \dot{y}^{\,2}
        +
        \dot{z}^{\,2}
    \right)
    \\[0.4em]
    &\quad
    -
    m\omega
    \Bigl[
        y
        \bigl(
            \cos(\phi_0)\dot{x}
            +
            \sin(\phi_0)\dot{z}
        \bigr)
        -
        \dot{y}
        \bigl(
            \cos(\phi_0)x
            +
            \sin(\phi_0)z
        \bigr)
    \Bigr]
    -
    mgz.
  \end{aligned}
\]
  \subsection{Introducing the Pendulum Constraint}

The equations derived so far describe the motion of a free particle in the rotating
frame of the Earth. To obtain the equations of motion of the Foucault pendulum,
we must now introduce the geometric constraint imposed by the rigid wire.

One possible approach would be to introduce spherical coordinates centred at the
point of suspension. In these coordinates the pendulum constraint is automatically
satisfied by fixing the radial coordinate equal to the length of the pendulum.
However, this choice is not particularly convenient. The equilibrium configuration
of the pendulum corresponds to the bob hanging vertically below the point of
suspension, which is precisely the pole of the spherical coordinate system. At this
point the azimuthal angle becomes undefined, making spherical coordinates poorly
suited for the linearization that follows.

Instead, we retain the local Cartesian coordinates $(x,y,z)$ introduced in the
previous section. We choose the origin of this coordinate system to coincide with
the equilibrium position of the pendulum bob. Consequently, the point of
suspension has coordinates
 
\[
(0,0,l),  
\]

where \(l\) denotes the length of the pendulum.

\begin{figure}[h] 
  \centering
  \includegraphics {pendulum.png}
  \caption{Local Cartesian coordinates of the pendulum.}
\end{figure}

The motion of the bob is therefore constrained to remain on the sphere of radius
\(l\) centred at the point of suspension. The corresponding holonomic constraint
is

\[
g(x,y,z)
=
x^2+y^2+(z-l)^2-l^2
=
0.
\]

We introduce a Lagrange multiplier \(\lambda(t)\) and define the augmented
Lagrangian

\[
F 
=
L 
-
\lambda\,g .
\]

The equations of motion for the generalized coordinates
\[
q_i=(x,y,z)
\]
are

\[
\frac{d}{dt}
\left(
\frac{\partial F}{\partial \dot q_i}
\right)
-
\frac{\partial F}{\partial q_i}
=
0.
\]

These equations are supplemented by the constraint

\[
g(x,y,z)=0.
\]

Applying the Euler--Lagrange equations to the augmented Lagrangian immediately yields

\begin{equation}
\left\{
\begin{aligned}
m\ddot{x}
-2m\omega\cos(\phi_0)\,\dot{y}
-2\lambda x
&=0,
\\[0.4em]
m\ddot{y}
+
2m\omega
\left(
\cos(\phi_0)\,\dot{x}
+
\sin(\phi_0)\,\dot{z}
\right)
-
2\lambda y
&=0,
\\[0.4em]
m\ddot{z}
-
2m\omega\sin(\phi_0)\,\dot{y}
+
mg
-
2\lambda(z-l)
&=0,
\\[0.4em]
x^2+y^2+(z-l)^2
&=
l^2.
\end{aligned}
\right.
\label{eq:foucault-constrained}
\end{equation}

\subsection{Gravity Only Case}

Before considering the complete system, it is useful to make a short detour and
examine the pendulum in the absence of the Earth's rotation. This simpler case
will clarify both the role of the Lagrange multiplier and the linearization of the
constraint around the equilibrium configuration.

Setting \(\omega=0\), the constrained equations of motion become

\begin{equation}
  \left\{
  \begin{aligned}
    m\ddot{x}-2\lambda x
    &=0,
    \\
    m\ddot{y}-2\lambda y
    &=0,
    \\
    m\ddot{z}+mg-2\lambda(z-l)
    &=0,
    \\
    x^2+y^2+(z-l)^2
    &=l^2.
  \end{aligned}
  \right.
  \label{eq:pendulum-gravity-constrained}
\end{equation}

The equilibrium position of the bob is

\[
  x=y=z=0.
\]

We therefore consider small oscillations around this configuration and write

\[
  z=\delta z,
  \qquad
  |\delta z|\ll l.
\]

Expanding the constraint gives

\[
  x^2+y^2+z^2-2lz=0.
\]

Since \(z\ll l\), the term \(z^2\) can be neglected, and the constraint reduces to

\begin{equation}
  x^2+y^2=2lz.
  \label{eq:pendulum-linearized-constraint}
\end{equation}


\begin{figure}[h] 
  \centering
  \includegraphics {secondOrderZ1.png}
  \caption{Why the vertical displacement is second order.}
\end{figure}


This relation shows that, if \(x\) and \(y\) are treated as first-order
displacements, then \(z\) is a second-order quantity:

\[
  z
  =
  \frac{x^2+y^2}{2l}.
\]

We now eliminate the Lagrange multiplier from the equations for \(y\) and \(z\).
Multiplying the equation for \(y\) by \(z-l\), multiplying the equation for \(z\)
by \(y\), and subtracting the latter from the former, we obtain

\[
  m\ddot{y}(z-l)
  -
  m\ddot{z}\,y
  -
  mgy
  =
  0.
\]

Equivalently,

\[
  -ml\ddot{y}
  +
  mz\ddot{y}
  -
  m\ddot{z}\,y
  -
  mgy
  =
  0.
\]

The displacement \(z\) and its derivatives are of second order in the horizontal
oscillation amplitude. Consequently, the terms \(z\ddot{y}\) and
\(\ddot{z}\,y\) are of higher order and can be neglected in the linear
approximation. We are therefore left with

\[
  -ml\ddot{y}-mgy=0,
\]

or

\begin{equation}
  \ddot{y}+\frac{g}{l}y=0.
  \label{eq:pendulum-y-equation}
\end{equation}

Repeating the same procedure with the equations for \(x\) and \(z\) gives

\begin{equation}
  \ddot{x}+\frac{g}{l}x=0.
  \label{eq:pendulum-x-equation}
\end{equation}

Thus, to first order in the horizontal displacements, the two components of the
motion satisfy the independent harmonic-oscillator equations

\begin{equation}
  \left\{
  \begin{aligned}
    \ddot{x}+\frac{g}{l}x&=0,
    \\
    \ddot{y}+\frac{g}{l}y&=0.
  \end{aligned}
  \right.
  \label{eq:ordinary-pendulum-linearized}
\end{equation}

The corresponding natural angular frequency is

\[
  \omega_0=\sqrt{\frac{g}{l}}.
\]

\subsection{Linearized Equations}

We now return to the complete equations of motion, including the Coriolis force,
and repeat the same procedure used in the gravity-only case.

We begin with the equations for \(y\) and \(z\),

\[
  m\ddot{y}
  +
  2m\omega
  \left(
    \cos\phi_0\,\dot{x}
    +
    \sin\phi_0\,\dot{z}
  \right)
  -
  2\lambda y
  =
  0,
\]

and

\[
  m\ddot{z}
  -
  2m\omega\sin\phi_0\,\dot{y}
  +
  mg
  -
  2\lambda(z-l)
  =
  0.
\]

Multiplying the first equation by \(z-l\), the second by \(y\), and
subtracting the latter from the former, we eliminate the Lagrange multiplier
and obtain

\[
  \begin{aligned}
    m\ddot{y}(z-l)
    &+
    2m\omega
    \left(
      \cos\phi_0\,\dot{x}
      +
      \sin\phi_0\,\dot{z}
    \right)(z-l)
    \\
    &-
    m\ddot{z}\,y
    +
    2m\omega\sin\phi_0\,\dot{y}\,y
    -
    mgy
    =
    0.
  \end{aligned}
\]

As in the gravity-only case, we consider small vertical displacements around
the equilibrium position. We therefore use

\[
  z-l\simeq-l
\]

and neglect \(z\), \(\dot z\), and \(\ddot z\), since they are of higher
order in the horizontal oscillation amplitude. The equation consequently
reduces to

\[
  -ml\ddot{y}
  -
  2m\omega l\cos\phi_0\,\dot{x}
  +
  2m\omega\sin\phi_0\,y\dot{y}
  -
  mgy
  =
  0.
\]

Dividing by \(-ml\), we obtain

\begin{equation}
  \ddot{y}
  +
  2\omega\cos\phi_0\,\dot{x}
  -
  \frac{2\omega}{l}\sin\phi_0\,y\dot{y}
  +
  \frac{g}{l}y
  =
  0.
  \label{eq:foucault-y-partially-linearized}
\end{equation}

We proceed in the same way with the equations for \(x\) and \(z\).
Multiplying the equation for \(x\) by \(z-l\), the equation for \(z\) by
\(x\), and subtracting, gives

\[
  \begin{aligned}
    m\ddot{x}(z-l)
    &-
    2m\omega\cos\phi_0\,\dot{y}(z-l)
    -
    m\ddot{z}\,x
    \\
    &+
    2m\omega\sin\phi_0\,x\dot{y}
    -
    mgx
    =
    0.
  \end{aligned}
\]

Using again \(z-l\simeq-l\) and neglecting \(\ddot z\), we find

\[
  -ml\ddot{x}
  +
  2m\omega l\cos\phi_0\,\dot{y}
  +
  2m\omega\sin\phi_0\,x\dot{y}
  -
  mgx
  =
  0,
\]

and therefore

\begin{equation}
  \ddot{x}
  -
  2\omega\cos\phi_0\,\dot{y}
  -
  \frac{2\omega}{l}\sin\phi_0\,x\dot{y}
  +
  \frac{g}{l}x
  =
  0.
  \label{eq:foucault-x-partially-linearized}
\end{equation}

The constraint must still be imposed. Its leading nonvanishing approximation
around the equilibrium position is

\[
  x^2+y^2+z^2-2lz=0
  \qquad\Longrightarrow\qquad
  2lz=x^2+y^2,
\]

where \(z^2\) has been neglected with respect to \(2lz\).

We therefore obtain the partially linearized constrained system

\begin{equation}
  \left\{
  \begin{aligned}
    \ddot{y}
    +
    2\omega\cos\phi_0\,\dot{x}
    -
    \frac{2\omega}{l}\sin\phi_0\,y\dot{y}
    +
    \frac{g}{l}y
    &=
    0,
    \\[0.5em]
    \ddot{x}
    -
    2\omega\cos\phi_0\,\dot{y}
    -
    \frac{2\omega}{l}\sin\phi_0\,x\dot{y}
    +
    \frac{g}{l}x
    &=
    0,
    \\[0.5em]
    2lz
    &=
    x^2+y^2.
  \end{aligned}
  \right.
  \label{eq:foucault-partially-linearized-system}
\end{equation}

We now complete the linearization of the equations of motion. Let the horizontal
displacements \(x\) and \(y\) be small compared with the length of the pendulum,

\[
  |x|\ll l,
  \qquad
  |y|\ll l.
\]

The velocities \(\dot{x}\) and \(\dot{y}\) are of the same perturbative order as
the corresponding displacements. Consequently, the mixed terms

\[
  x\dot{y},
  \qquad
  y\dot{y}
\]

are of second order in the oscillation amplitude and can be neglected in the
linear approximation.

The equations of motion then reduce to

\begin{equation}
  \left\{
  \begin{aligned}
    \ddot{y}
    +
    2\omega\cos\phi_0\,\dot{x}
    +
    \frac{g}{l}y
    &=
    0,
    \\[0.4em]
    \ddot{x}
    -
    2\omega\cos\phi_0\,\dot{y}
    +
    \frac{g}{l}x
    &=
    0,
    \\[0.4em]
    2lz
    &=
    x^2+y^2.
  \end{aligned}
  \right.
  \label{eq:foucault-linearized-equations}
\end{equation}

\subsection{Solution of the Equations of Motion}

We transform the second-order system into an equivalent first-order system by
introducing the auxiliary variables

\[
u=\dot{x},
\qquad
v=\dot{y}.
\]


The vertical constraint no longer enters explicitly into the horizontal equations
of motion and will therefore be omitted in what follows. It can always be
recovered afterwards from

\[
2lz=x^2+y^2.
\]
Defining

\[
\alpha=2\omega\cos\phi_0,
\qquad
\beta=\frac{g}{l},
\]

the equations become

\[
\left\{
\begin{aligned}
\dot{x} &= u,\\
\dot{y} &= v,\\
\dot{u} &= \alpha v-\beta x,\\
\dot{v} &= -\alpha u-\beta y.
\end{aligned}
\right.
\]

The system can now be written in matrix form as

\[
\frac{d}{dt}
\begin{pmatrix}
x\\
y\\
u\\
v
\end{pmatrix}
=
\begin{pmatrix}
0 & 0 & 1 & 0\\
0 & 0 & 0 & 1\\
-\beta & 0 & 0 & \alpha\\
0 & -\beta & -\alpha & 0
\end{pmatrix}
\begin{pmatrix}
x\\
y\\
u\\
v
\end{pmatrix}.
\]

The general solution of this system is determined by the eigenvalues and
eigenvectors of the coefficient matrix. A straightforward calculation shows
that the four eigenvalues are

\begin{align}
  \gamma_{1,2} &= \pm i\omega_1, \\
  \gamma_{3,4} &= \pm i\omega_2,
\end{align}

where

\begin{align}
  \omega_1
  &=
  \sqrt{
    \frac{
      \alpha^2
      +
      2\beta
      +
      \alpha\sqrt{\alpha^2+4\beta}
    }{2}
  },
  \\
  \omega_2
  &=
  \sqrt{
    \frac{
      \alpha^2
      +
      2\beta
      -
      \alpha\sqrt{\alpha^2+4\beta}
    }{2}
  }.
\end{align}

The system therefore possesses four distinct eigenvalues corresponding to two
independent oscillation frequencies.

It is worth observing that, in the absence of the Coriolis force
(\(\alpha=0\)), the two frequencies become equal,

\[
\omega_1=\omega_2=\sqrt{\beta}=\sqrt{\frac{g}{l}},
\]

and the eigenvalues become doubly degenerate. The Coriolis force therefore
removes this degeneracy, splitting the single natural frequency of the
ordinary pendulum into two distinct oscillation frequencies.\par


\begin{figure}[h] 
  \centering
  \includegraphics {frequenciesSplitting1.png}
  \caption{Splitting of the normal-mode frequencies by the Coriolis force.}
\end{figure}

We now determine the corresponding eigenvectors. Since the procedure is
identical for each eigenvalue, we derive only the first one in detail.

For the eigenvalue

\[
\gamma_1=i\omega_1,
\]

the eigenvalue equation

\[
(A-\gamma_1 I)\vb{v}_1=0
\]

reads

\[
\begin{pmatrix}
i\omega_1 & 0 & -1 & 0 \\
0 & i\omega_1 & 0 & -1 \\
\beta & 0 & i\omega_1 & -\alpha \\
0 & \beta & \alpha & i\omega_1
\end{pmatrix}
\begin{pmatrix}
x\\
y\\
u\\
v
\end{pmatrix}
=
\begin{pmatrix}
0\\
0\\
0\\
0
\end{pmatrix}.
\]

The first two equations immediately give

\[
u=i\omega_1x,
\qquad
v=i\omega_1y.
\]

Substituting these expressions into the third equation yields

\[
(\beta-\omega_1^2)x
=
i\alpha\omega_1\,y,
\]

or

\[
y
=
i\,
\frac{\omega_1^2-\beta}{\alpha\omega_1}
\,x.
\]

Since eigenvectors are defined up to an arbitrary multiplicative constant, we
choose

\[
x=1.
\]

The corresponding eigenvector is therefore

\[
 
\vb{v}_1
=
\begin{pmatrix}
1\\[0.8em]
i\dfrac{\omega_1^2-\beta}{\alpha\omega_1}\\[1.0em]
i\omega_1\\[0.8em]
\dfrac{\beta-\omega_1^2}{\alpha}
\end{pmatrix}
 
\]
For convenience, we introduce the quantities

\[
\eta_1
=
\frac{\omega_1^2-\beta}{\alpha\omega_1},
\qquad
\eta_2
=
\frac{\omega_2^2-\beta}{\alpha\omega_2}.
\]

The corresponding eigenvectors are

\[
\vb{v}_1=
\begin{pmatrix}
1\\
i\eta_1\\
i\omega_1\\
-\omega_1\eta_1
\end{pmatrix},
\qquad
\vb{v}_3=
\begin{pmatrix}
1\\
i\eta_2\\
i\omega_2\\
-\omega_2\eta_2
\end{pmatrix},
\]

while the remaining two eigenvectors are simply their complex conjugates,

\[
\vb{v}_2=\overline{\vb{v}}_1,
\qquad
\vb{v}_4=\overline{\vb{v}}_3.
\]

The general solution of the system can therefore be written as

\[
\begin{pmatrix}
x\\
y\\
u\\
v
\end{pmatrix}
=
c_1\vb{v}_1e^{i\omega_1t}
+
c_2\vb{v}_2e^{-i\omega_1t}
+
c_3\vb{v}_3e^{i\omega_2t}
+
c_4\vb{v}_4e^{-i\omega_2t},
\]

where the constants \(c_1,c_2,c_3,c_4\) are determined by the initial
conditions.

To illustrate the motion, we consider the pendulum released from rest after a
small displacement \(d\) along the \(x\)-axis,

\[
\begin{pmatrix}
x(0)\\
y(0)\\
u(0)\\
v(0)
\end{pmatrix}
=
\begin{pmatrix}
d\\
0\\
0\\
0
\end{pmatrix}
\]

This choice is not restrictive. Since the equations of motion are invariant
under rotations in the horizontal plane, any other initial displacement can be
obtained from this solution by a suitable rotation of the coordinate system.

Imposing these initial conditions uniquely determines the coefficients
\(c_1,c_2,c_3,c_4\). After straightforward algebra, the horizontal motion is
found to be

\begin{align}
x(t)
&=
\frac{d}{2}
\left[
\frac{\alpha}{\sqrt{\alpha^2+4\beta}}
\left(
\cos(\omega_2t)
-
\cos(\omega_1t)
\right)
+
\cos(\omega_1t)
+
\cos(\omega_2t)
\right],
\\[0.8em]
y(t)
&=
\frac{d}{\sqrt{\alpha^2+4\beta}}
\left[
\omega_1\sin(\omega_2t)
-
\omega_2\sin(\omega_1t)
\right].
\end{align}
 
The expressions obtained above can be rewritten in a more revealing form, in
which the half-sum and half-difference of the two eigenfrequencies appear
explicitly.

We begin with the expression for \(y(t)\),

\[
  y(t)
  =
  \frac{d}{\sqrt{\alpha^2+4\beta}}
  \left[
    \omega_1\sin(\omega_2t)
    -
    \omega_2\sin(\omega_1t)
  \right].
\]

The quantity inside the brackets can be rearranged as

\[
  \begin{aligned}
    \omega_1\sin(\omega_2t)
    -
    \omega_2\sin(\omega_1t)
    &=
    \frac{\omega_1-\omega_2}{2}
    \left[
      \sin(\omega_1t)
      +
      \sin(\omega_2t)
    \right]
    \\
    &\quad
    -
    \frac{\omega_1+\omega_2}{2}
    \left[
      \sin(\omega_1t)
      -
      \sin(\omega_2t)
    \right].
  \end{aligned}
\]

We now use the sum-to-product identities

\[
  \sin a+\sin b
  =
  2\sin\left(\frac{a+b}{2}\right)
  \cos\left(\frac{a-b}{2}\right),
\]

and

\[
  \sin a-\sin b
  =
  2\cos\left(\frac{a+b}{2}\right)
  \sin\left(\frac{a-b}{2}\right).
\]

It follows that

\[
  \begin{aligned}
    y(t)
    &=
    \frac{d}{\sqrt{\alpha^2+4\beta}}
    \Biggl[
      (\omega_1-\omega_2)
      \cos\left(
        \frac{\omega_1-\omega_2}{2}\,t
      \right)
      \sin\left(
        \frac{\omega_1+\omega_2}{2}\,t
      \right)
      \\
      &\hspace{7em}
      -
      (\omega_1+\omega_2)
      \cos\left(
        \frac{\omega_1+\omega_2}{2}\,t
      \right)
      \sin\left(
        \frac{\omega_1-\omega_2}{2}\,t
      \right)
    \Biggr].
  \end{aligned}
\]

We proceed in the same way with \(x(t)\),

\[
  x(t)
  =
  \frac{d}{2}
  \left[
    \frac{\alpha}{\sqrt{\alpha^2+4\beta}}
    \left(
      \cos(\omega_2t)-\cos(\omega_1t)
    \right)
    +
    \cos(\omega_1t)+\cos(\omega_2t)
  \right].
\]

Using the identities

\[
  \cos a+\cos b
  =
  2\cos\left(\frac{a+b}{2}\right)
  \cos\left(\frac{a-b}{2}\right),
\]

and

\[
  \cos b-\cos a
  =
  2\sin\left(\frac{a+b}{2}\right)
  \sin\left(\frac{a-b}{2}\right),
\]

we find

\[
  \begin{aligned}
    x(t)
    &=
    d
    \Biggl[
      \cos\left(
        \frac{\omega_1-\omega_2}{2}\,t
      \right)
      \cos\left(
        \frac{\omega_1+\omega_2}{2}\,t
      \right)
      \\
      &\hspace{5em}
      +
      \frac{\alpha}{\sqrt{\alpha^2+4\beta}}
      \sin\left(
        \frac{\omega_1-\omega_2}{2}\,t
      \right)
      \sin\left(
        \frac{\omega_1+\omega_2}{2}\,t
      \right)
    \Biggr].
  \end{aligned}
\]

The solution can therefore be written in the final form

\begin{equation}
  \left\{
  \begin{aligned}
    x(t)
    &=
    d
    \Biggl[
      \cos\left(
        \frac{\omega_1-\omega_2}{2}\,t
      \right)
      \cos\left(
        \frac{\omega_1+\omega_2}{2}\,t
      \right)
      \\
      &\hspace{5em}
      +
      \frac{\alpha}{\sqrt{\alpha^2+4\beta}}
      \sin\left(
        \frac{\omega_1-\omega_2}{2}\,t
      \right)
      \sin\left(
        \frac{\omega_1+\omega_2}{2}\,t
      \right)
    \Biggr],
    \\[0.8em]
    y(t)
    &=
    \frac{d}{\sqrt{\alpha^2+4\beta}}
    \Biggl[
      (\omega_1-\omega_2)
      \cos\left(
        \frac{\omega_1-\omega_2}{2}\,t
      \right)
      \sin\left(
        \frac{\omega_1+\omega_2}{2}\,t
      \right)
      \\
      &\hspace{5em}
      -
      (\omega_1+\omega_2)
      \cos\left(
        \frac{\omega_1+\omega_2}{2}\,t
      \right)
      \sin\left(
        \frac{\omega_1-\omega_2}{2}\,t
      \right)
    \Biggr].
  \end{aligned}
  \right.
  \label{eq:foucault-sum-difference-solution}
\end{equation}

We now examine the behavior of the two eigenfrequencies when the Coriolis
coupling is small. Since

\[
  \alpha=2\omega\cos\phi_0
\]

is much smaller than the natural frequency of the pendulum, we expand
\(\omega_1\) and \(\omega_2\) around \(\alpha=0\).

At \(\alpha=0\), both frequencies reduce to

\[
  \omega_1(0)
  =
  \omega_2(0)
  =
  \sqrt{\beta}.
\]

For the first frequency,

\[
  \omega_1
  =
  \sqrt{
    \frac{
      \alpha^2
      +
      2\beta
      +
      \alpha\sqrt{\alpha^2+4\beta}
    }{2}
  },
\]

and its derivative at \(\alpha=0\) is

\[
  \left.
  \frac{d\omega_1}{d\alpha}
  \right|_{\alpha=0}
  =
  \frac{1}{2}.
\]

Similarly,

\[
  \omega_2
  =
  \sqrt{
    \frac{
      \alpha^2
      +
      2\beta
      -
      \alpha\sqrt{\alpha^2+4\beta}
    }{2}
  },
\]

and

\[
  \left.
  \frac{d\omega_2}{d\alpha}
  \right|_{\alpha=0}
  =
  -\frac{1}{2}.
\]

The Taylor expansions to first order in \(\alpha\) are therefore

\begin{align}
  \omega_1
  &=
  \sqrt{\beta}
  +
  \frac{\alpha}{2}
  +
  O(\alpha^2),
  \\
  \omega_2
  &=
  \sqrt{\beta}
  -
  \frac{\alpha}{2}
  +
  O(\alpha^2).
\end{align}

It follows immediately that

\begin{align}
  \omega_1+\omega_2
  &=
  2\sqrt{\beta}
  +
  O(\alpha^2),
  \\
  \omega_1-\omega_2
  &=
  \alpha
  +
  O(\alpha^3).
\end{align}

Thus, to the required order,

\[
  \frac{\omega_1+\omega_2}{2}
  \simeq
  \sqrt{\beta},
  \qquad
  \frac{\omega_1-\omega_2}{2}
  \simeq
  \frac{\alpha}{2}.
\]

Substituting these expressions into
\eqref{eq:foucault-sum-difference-solution}, we obtain

\begin{equation}
  \left\{
  \begin{aligned}
    x(t)
    &\simeq
    d
    \Biggl[
      \cos\left(\frac{\alpha}{2}t\right)
      \cos\left(\sqrt{\beta}\,t\right)
      +
      \frac{\alpha}{\sqrt{\alpha^2+4\beta}}\,
      \sin\left(\frac{\alpha}{2}t\right)
      \sin\left(\sqrt{\beta}\,t\right)
    \Biggr],
    \\[0.8em]
    y(t)
    &\simeq
    \frac{d}{\sqrt{\alpha^2+4\beta}}
    \Biggl[
      \alpha
      \cos\left(\frac{\alpha}{2}t\right)
      \sin\left(\sqrt{\beta}\,t\right)
      \\
      &\hspace{9em}
      -
      2\sqrt{\beta}\,
      \cos\left(\sqrt{\beta}\,t\right)
      \sin\left(\frac{\alpha}{2}t\right)
    \Biggr].
  \end{aligned}
  \right.
  \label{eq:foucault-frequency-expanded-solution}
\end{equation}

We now perform the final simplification. Since

\[
  \frac{\alpha}{\sqrt{\alpha^2+4\beta}}
  \simeq
  \frac{\alpha}{2\sqrt{\beta}}
  \ll 1,
\]

the second term in \(x(t)\) is a small correction to its leading term and can
be neglected. Similarly, in the expression for \(y(t)\), the term proportional
to \(\alpha\) is negligible compared with the term proportional to
\(2\sqrt{\beta}\). We also use

\[
  \sqrt{\alpha^2+4\beta}
  \simeq
  2\sqrt{\beta}.
\]

It is important, however, not to expand the functions

\[
  \cos\left(\frac{\alpha}{2}t\right),
  \qquad
  \sin\left(\frac{\alpha}{2}t\right),
\]

because their variation occurs over the long time scale associated with the
precession of the pendulum.

The horizontal motion therefore reduces to

\begin{equation}
  \left\{
  \begin{aligned}
    x(t)
    &=
    d
    \cos\left(\frac{\alpha}{2}t\right)
    \cos\left(\sqrt{\beta}\,t\right),
    \\[0.5em]
    y(t)
    &=
    -d
    \sin\left(\frac{\alpha}{2}t\right)
    \cos\left(\sqrt{\beta}\,t\right).
  \end{aligned}
  \right.
  \label{eq:foucault-final-solution-alpha-beta}
\end{equation}

Recalling that

\[
  \alpha=2\omega\cos\phi_0,
  \qquad
  \beta=\frac{g}{l},
\]

we finally obtain

\begin{equation}
  \left\{
  \begin{aligned}
    x(t)
    &=
    d
    \cos\left(\omega\cos\phi_0\,t\right)
    \cos\left(\sqrt{\frac{g}{l}}\,t\right),
    \\[0.5em]
    y(t)
    &=
    -d
    \sin\left(\omega\cos\phi_0\,t\right)
    \cos\left(\sqrt{\frac{g}{l}}\,t\right).
  \end{aligned}
  \right.
  \label{eq:foucault-final-solution}
\end{equation}

\subsection{Precession of the Oscillation Direction}

The form of the solution admits a particularly simple geometric
interpretation. Introducing the rotation matrix

\[
R(\vartheta)
=
\begin{pmatrix}
\cos\vartheta & -\sin\vartheta\\
\sin\vartheta & \cos\vartheta
\end{pmatrix},
\]

the horizontal motion can be written as

\begin{equation}
  \begin{pmatrix}
    x(t)\\
    y(t)
  \end{pmatrix}
  =
  R\left(-\omega\cos\phi_0\,t\right)
  \begin{pmatrix}
    d\cos\left(\sqrt{\frac{g}{l}}\,t\right)\\
    0
  \end{pmatrix}.
  \label{eq:foucault-rotating-direction}
\end{equation}

The pendulum therefore performs a rapid oscillation with angular frequency

\[
\sqrt{\frac{g}{l}},
\]

along the horizontal unit vector

\[
\vb n(t)
=
\cos\left(\omega\cos\phi_0\,t\right)\vb e_x
-
\sin\left(\omega\cos\phi_0\,t\right)\vb e_y.
\]

This direction rotates slowly with angular velocity

\[
-\omega\cos\phi_0.
\]

Together with the local vertical direction, \(\vb n(t)\) determines the
instantaneous plane of oscillation of the pendulum. It is in this precise
sense that the plane of oscillation is said to rotate relative to the Earth.

The underlying dynamical mechanism, however, is the splitting of the two
circular normal modes by the Coriolis force. In the absence of the Earth's
rotation, the two modes have the same frequency. The Coriolis coupling removes
this degeneracy, and the gradual accumulation of their relative phase produces
the slow precession of the oscillation direction.


\end{document}
